\begin{figure}[t]
\centering
\begin{tikzpicture}[font=\small,
  floor/.style={draw,rounded corners=2pt,minimum width=1.5cm,minimum height=0.6cm},
  coord/.style={draw,minimum width=0.55cm,minimum height=0.45cm,inner sep=1pt},
  pass/.style={coord,fill=gray!55},
  fail/.style={coord,pattern=north east lines,pattern color=black!70},
  skip/.style={coord,draw=gray!60,densely dotted},
  >=Stealth]
  \foreach \k/\y in {300/0,200/-1,150/-2,100/-3,75/-4} {
    \node[floor] (f\k) at (0,\y) {$k=\k$};
  }
  % components: single, pooled-history, full
  \foreach \k/\y/\a/\b/\c in {300/0/pass/pass/pass,200/-1/pass/pass/pass,150/-2/pass/pass/pass,100/-3/pass/pass/fail,75/-4/skip/skip/skip} {
    \node[\a] at (1.8,\y) {};
    \node[\b] at (3.0,\y) {};
    \node[\c] at (4.2,\y) {};
  }
  \node[align=center,font=\footnotesize] at (1.8,0.9) {one\\release};
  \node[align=center,font=\footnotesize] at (3.0,0.9) {pooled,\\history};
  \node[align=center,font=\footnotesize] at (4.2,0.9) {plus\\readings};
  \foreach \a/\b in {300/200,200/150,150/100,100/75} {\draw[->] (f\a) -- (f\b);}
  \node[anchor=west,align=left] at (5.0,-2) {last floor that passed:\\ chosen floor $k=150$};
  \node[anchor=west,align=left] at (5.0,-3) {one attack fails: chain stops};
  \node[anchor=west,align=left] at (5.0,-4) {never opened};
  % legend
  \node[pass] at (-0.3,-5.3) {}; \node[anchor=west] at (0.05,-5.3) {certified};
  \node[fail] at (2.0,-5.3) {}; \node[anchor=west] at (2.35,-5.3) {not certified within the cap};
  \node[skip] at (7.6,-5.3) {}; \node[anchor=west] at (7.95,-5.3) {not audited};
\end{tikzpicture}
\caption{A release chain chooses the cell-count floor. Floors are audited from the
strictest down; at each floor every declared attack is a coordinate audited on
fresh targets at the full error level, and the floor passes only if all of them
certify. The chain stops at the first floor that fails and chooses the last one
that passed (Algorithm~\ref{alg:chain}). The pattern shown is the one most streams
followed in the credit-register audit of the companion article.}
\label{fig:chain}
\end{figure}
